
M://Ma come funziona se taglia in 2 posti?
M:suppongo res ritorni il polygono tagliato, ah ma comunque ce ne sarebbero 2
S:in effetti sono molto triggerato

M:aspè proviamo a pseudo coddarlo
S:Paxxo

P:Raga sono frocio

M: No dai non dire così, la tua capacità di scrivere codice intelligibile, non dovrebbe influire sulla tua sessualità
S:Ma che state dicendo? Non capisco se siete seri o no
M: credo lui fosse molto serio

hai un vettore di punti
per ogni punto del poligono ti salvi quello corrente e ciclicamente il precedente
che fa cross? ah il prodotto vettoriale? //Adesso ce guardo
Cioè non è proprio il prodotto vettoriale, penso che veda tipo l'orientamento (non sessuale)
Ok quindi se l'orientamento tra cur e prev cambia, allora faccio:
    aggiunge l'intersezione 
altrimenti:
    aggiungo il punto corrente

no ma che sarebbe intersection? AHHH è per reference (per quello non capivo)
ah praticamente prende solo il poligono che ha orientamento positivo

ok quindi per avere l'altro basta mettere il ! prima di side nel secondo if

ok ma perché fare una funzione così merdosa?
ma tra l'altro non funziona solo se è convesso?

sì

basta fare una convex decomposition

Ma sta roba ce la lasciamo sulla libreria? finche non ne troviamo una migliore

no ma io intendevo sti commenti omofobi

sì, quelli che stamperanno la libreria devono sapere che ci teniamo alla nostra sessualità e che non devono omosessualizzarci

ma infatti alle swerc scorse hanno distribuito le magliette VIOLA, quest'anno le cose saranno diverse

era viola perché siamo autistici come la persona da cui questo colore prende il nome

lel

comunque non so se ce l'abbiamo, ma forse scomporre i poligoni in poligoni convessi è una funzione utile

mhhhhhhhhhhhhhhhhhhhhhhhhhhhhhhhh, intendi triangoli?

ah cavolo è vero, è ancora più facile

allora basta triangolarizzare il poligono, poi usare sta funzione, sommare i perimetri e togliere (tutti i lati aggiunti dei triangoli e i tagli fatti)*2

Ma questo risolve il problema, non adatta la funzione a poligoni generici

ma questa funzione cosa vuoi che faccia? ritornare tutti i nuovi poligoni?

s ì
però ammetto che sia molto più cancra così




triangolizzi, tagli e incolli. I lati dei triangoli aggiunti sono speciali e dopo il taglio diventano "appiccicosi". è un po' un casino perchè si dividono in 2

è shi è troppo cancro da avere in libreria me sa? (e fu così che alle swerc venne fuori questo esatto problema)

usi questa funzione e poi i segmenti il cui punto medio ricadono fuori dal poligono, li cancelli

ma magari puoi semplicemente chiudere un poligono quando crossa

anche, sì hai ragione

c'è sempre il problema dei punti che cadono sulla retta? (e non dirmi di aggiungere noise pliz) forse non è un problema in realtà

non so come funzioni cross, ma se trova quelli sulla retta, allora basta chiuderli come gli altri

Overleaf >> Zoom! (no seriamente semmai la scriviamo con l'ausilio di discord) 

Ma in realtà ora dovrei andare

Ma sì un altro giorno, ovviamente tutto il discorso rimane

Ah comunque non funziona lol
\fi